\subsection{The Hahn-Banach theorem and its applications}

THEOREM 1 (Hahn-Banach). --- Let V be a vector subspace of E, a complex vector space, and let f be a (complex) linear form on V and p a semi-norm on E such that \(|f(y)| \leqslant p(y)\) for all \(y \in V\) . Then there exists a linear form \(f_{1}\) on E extending f and such that \(|f_{1}(x)| \leqslant p(x)\) for all \(x \in E\) .

For \(g = \mathcal{R}f\) is a real linear form defined in \(V\) and satisfying \(|g(y)| \leqslant p(y)\) at every point of \(V\) ; therefore there exists a real linear form \(g_1\) in \(E\) extending \(g\) and such that \(|g_1(x)| \leqslant p(x)\) for all \(x \in E\) (II, p.~23, cor. 1). Let \(f_1(x) = g_1(x) - ig_1(ix)\) be the complex linear form on \(E\) of which \(g_1\) is the real part (II, p.~61). For all real \(\vartheta\)

\[
\left| \mathcal {R} (e ^ {i \vartheta} f _ {1} (x)) \right| = \left| \mathcal {R} (f _ {1} (e ^ {i \vartheta} x)) \right| = \left| g _ {1} (e ^ {i \vartheta} x) \right| \leqslant p (e ^ {i \vartheta} x) = p (x)
\]

since \(p\) is a semi-norm on the complex space \(\mathbf{E}\) ; this implies the relation \(|f_1(x)| \leqslant p(x)\) , and the theorem is proved.

COROLLARY 1. --- Let \(x_0\) be a point of a complex topological vector space \(E\) and \(p\) be a continuous semi-norm in \(E\) ; then there exists a continuous (complex) linear form \(f\) defined in \(E\) , such that \(f(x_0) = p(x_0)\) and \(|f(x)| \leqslant p(x)\) for all \(x \in E\) .

COROLLARY 2. --- Let V be a vector subspace of a complex locally convex space E and f be a (complex) linear form defined and continuous in V; then there exists a continuous linear form \(f_{1}\) defined in E and extending f.~If E is normed there exists such a form \(f_{1}\) that also satisfies \(\|f_{1}\|=\|f\|\) .

COROLLARY 3. --- Let M be a finite dimensional vector subspace of a Hausdorff complex locally convex space E. Then there exists a closed vector subspace N of E that is a topological complement of M in E.

The proofs using theorem 1, p.~24 are the same as those of II, p.~23, cor. 2 and cor. 3, p.~24, prop. 2 and p.~25, cor. 2.

PROPOSITION 1. --- Let A be an open non-empty convex set in a complex topological vector space E and M be a non-empty (complex) linear variety that does not meet A. Then there exists a closed complex hyperplane H that contains M and does not meet A.

We can suppose that \(0 \in M\) . Then there exists a closed real hyperplane \(H_{0}\) containing M and not meeting A (II, p.~36; th. 1). As M = iM, the closed complex hyperplane \(\mathbf{H} = \mathbf{H}_{0} \cap (i\mathbf{H}_{0})\) has the properties required.

COROLLARY. --- In a complex locally convex space E, every closed complex linear variety M is the intersection of the closed complex hyperplanes which contain it.

In fact, for all \(x \notin M\) , there exists a convex open neighbourhood V of x that does not meet M, and thus there exists a closed complex hyperplane H containing M and not meeting V; a fortiori H does not contain x.

PROPOSITION 2. --- Let A be a non-empty balanced open convex set of a complex topological vector space E, and B be a non-empty convex set that does not meet A. Then there exists a continuous complex linear form f on E and a number \(\alpha > 0\) such that \(|f(x)| < \alpha\) in A and \(|f(y)| \geqslant \alpha\) in B.

For, there exists a continuous real linear form \(g\) on \(\mathbf{E}\) and a real number \(\alpha\) such that \(g(x) < \alpha\) in \(\mathbf{A}\) and \(g(y) \geqslant \alpha\) in \(\mathbf{B}\) (II, p.~37, prop. 1). As \(0 \in \mathbf{A}\) , we have \(\alpha > 0\) . We show that the continuous complex linear form \(f(x) = g(x) - ig(ix)\) and the number \(\alpha\) have the properties required. For, since \(\mathcal{R}f = g\) , we have \(|f(y)| \geqslant \alpha\) in \(\mathbf{B}\) . On the other hand, for all \(x \in \mathbf{A}\) and all real \(\vartheta\) , the point \(e^{i\vartheta}x\) belongs to \(\mathbf{A}\) , since \(\mathbf{A}\) is balanced, and we have \(f(x) = e^{-i\vartheta}f(e^{i\vartheta}x)\) ; then there exists a number \(\vartheta\) such that \(|f(x)| = \mathcal{R}(e^{i\vartheta}f(x)) = g(e^{i\vartheta}x) < \alpha\) , and the proposition follows.

PROPOSITION 3. --- Let A be a balanced, closed, convex set in a complex locally convex space E and let K be a non-empty compact convex set in E that does not meet A. Then there exists a continuous complex linear form f on E and a number \(\alpha > 0\) such that \(|f(x)| < \alpha\) in A and \(|f(y)| > \alpha\) in K.

The proposition follows from II, p.~38, prop. 4 as prop. 2 follows from II, p.~37, prop. 1.

